% original PDF p. 47
\numpara{7.4}\label{XXIV.7.4}\textbf{An example}
As an example, we shall determine
\[
\underline{\Hom}_{S\text{-gr.}}(\SL_{2,S},\SL_{2,S}).
\]
\numpara{7.4.1}\label{XXIV.7.4.1}
Recall (Exp.~XX, no. 5) that $\SL_{2,S}$ is the $S$-group scheme consisting of matrices $\left(\begin{smallmatrix}a&b\\c&d\end{smallmatrix}\right)$ over $S$ satisfying $ad-bc=1$. A maximal torus $T$ is defined by the monomorphism $\alpha^*:\mathbf G_{\mathrm m,S}\to\SL_{2,S}$
\[
\alpha^*(z)=\begin{pmatrix}z&0\\0&z^{-1}\end{pmatrix}.
\]
A root of $G$ relative to $T$ is defined by $\alpha(\alpha^*(z))=z^2$, a corresponding monomorphism
\[
p:\mathbf G_{\mathrm a,S}\to\SL_{2,S}
\]
being defined by
\[
p(x)=\begin{pmatrix}1&x\\0&1\end{pmatrix}.
\]
Finally, the representative of the Weyl group corresponding to $u=p(1)$ is
\[
w=\begin{pmatrix}0&1\\-1&0\end{pmatrix}.
\]
Recall moreover (Exp.~XX, 6.2) that $\SL_{2,S}$ is generated by $T$, $p(\mathbf G_{\mathrm a,S})$ and $w$, subject to the relations
\[
\begin{aligned}
\alpha^*(z)p(x)\alpha^*(z^{-1})&=p(z^2x),\\
w\alpha^*(z)w^{-1}&=\alpha^*(z^{-1}),\\
w^2&=\alpha^*(-1),\\
(wu)^3&=1.
\end{aligned}
\]
\numpara{7.4.2}\label{XXIV.7.4.2}
Let $f$ be an endomorphism of $\SL_{2,S}$. It first defines a homomorphism $f\circ\alpha^*:\mathbf G_{\mathrm m,S}\to\SL_{2,S}$. The kernel $\Ker(f\circ\alpha^*)$ is a closed subgroup of $\mathbf G_{\mathrm m,S}$, hence is locally on $S$ equal to a $\bmu_{n,S}$ ($n\geq1$) or to $\mathbf G_{\mathrm m,S}$. After restricting $S$, one may therefore suppose that there exist an $n\in\NN$ and a monomorphism
\[
f':\mathbf G_{\mathrm m,S}\to\SL_{2,S}
\]
such that $f\circ\alpha^*(z)=f'(z^n)$.

By conjugacy of the monomorphisms $\mathbf G_{\mathrm m,S}\to\SL_{2,S}$, one can, after an extension of the base covering for the étale topology, find a section $g$ of $G$ such that $f\circ\alpha^*(z)=\mathrm{int}(g)\circ\alpha^*(z^n)$. Transforming $f$ by $\mathrm{int}(g)$, one is therefore reduced to the case where there exists an $n\in\NN$ such that $f\circ\alpha^*(z)=\alpha^*(z^n)$.
% original PDF p. 48
\numpara{7.4.3}\label{XXIV.7.4.3}
Consider now the morphism
\[
f\circ p:\mathbf G_{\mathrm a,S}\to\SL_{2,S}.
\]
It satisfies the condition
\[
\alpha^*(z^n)(f\circ p)(x)\alpha^*(z^n)^{-1}=(f\circ p)(z^2x).
\]
If $n=0$, it follows at once that $f\circ p$ is invariant under homotheties of $\mathbf G_{\mathrm a,S}$, hence constant. If $n\ne0$, one can apply Exp.~XXII, 4.1.9; $x\mto x^n$ is an endomorphism of $\mathbf G_{\mathrm a,S}$, and there exists a $\lambda\in\mathbf G_{\mathrm m,S}$ such that
\[
f\circ p(x)=p(\lambda x^n);
\]
this relation is moreover valid for $n=0$, taking $\lambda=1$. After again extending $S$, one can find a $z\in\mathbf G_{\mathrm m,S}$ such that $z^2=\lambda$. Replacing $f$ by $\mathrm{int}(\alpha^*(z))\circ f$, one is therefore reduced to the case where one has
\[
f\circ\alpha^*(z)=\alpha^*(z^n),\qquad f\circ p(x)=p(x^n).
\]
% typo?: source's n=0 assertion and direction of conjugation retained.
\numpara{7.4.4}\label{XXIV.7.4.4}
Finally, one must have $f(w)\alpha^*(z^n)f(w)^{-1}=\alpha^*(z^n)^{-1}$ and $(f(w)u)^3=1$. By Exp.~XX, 3.8, this implies $f(w)=w^n$. Since $G$ is generated by $T$, $p(\mathbf G_{\mathrm a,S})$ and $w$, this implies that for every $\left(\begin{smallmatrix}a&b\\c&d\end{smallmatrix}\right)\in G(S')$, $S'\to S$, one has
\[
f\begin{pmatrix}a&b\\c&d\end{pmatrix}=\begin{pmatrix}a^n&b^n\\c^n&d^n\end{pmatrix}.
\]
\numpara{7.4.5}\label{XXIV.7.4.5}
Summarizing the preceding discussion, one sees that locally on $S$ for the étale topology, one can find for every endomorphism $f$ of $\SL_{2,S}$ an (inner) automorphism $\phi$ of $\SL_{2,S}$, and an integer $n\geq0$ such that $f=\phi\circ F_n$, where
\[
F_n\begin{pmatrix}a&b\\c&d\end{pmatrix}=\begin{pmatrix}a^n&b^n\\c^n&d^n\end{pmatrix}.
\]
Observe that if $f=\phi\circ F_n$, the integer $n$ is determined by a fiber $f_s$ of $f$, for example by $\Ker(f_s)$. It follows that $n$ is a locally constant function on $S$.
% typo?: entrywise n-th power formula for n=0 retained as printed.
\numpara{7.4.6}\label{XXIV.7.4.6}
One deduces immediately that if $f$ is an endomorphism of $\SL_{2,S}$, then $S$ decomposes canonically as a sum of open and closed subschemes $S_0$, $S_1$, $S_{p^n}$ (where $p^n$ ranges over the powers $>0$ of prime numbers) such that:
\begin{enumeratei}
\item $f_{S_0}$ is the zero morphism,
\item $f_{S_1}$ is an isomorphism (= an inner automorphism),
\item $S_{p^n}$ is of characteristic $p$ and $f_{S_{p^n}}$ decomposes uniquely in the form $\phi\circ F_p^n$, where $\phi$ is an inner automorphism and $F_p$ the Frobenius endomorphism of $\SL_{2,\FF_p}$.
\end{enumeratei}
\numpara{7.4.7}\label{XXIV.7.4.7}
In other words, $\underline{\Hom}_{\ZZ\text{-gr.}}(\SL_{2,\ZZ},\SL_{2,\ZZ})$ is the sum scheme:
\begin{enumeratei}
\item of a scheme isomorphic to $\Spec(\ZZ)$,
\item of a scheme isomorphic to $\underline{\Aut}_{\ZZ\text{-gr.}}(\SL_{2,\ZZ})\simeq\ad(\SL_{2,\ZZ})$,
% original PDF p. 49
\item for each prime number $p$ and each integer $n>0$, of a scheme isomorphic to $\underline{\Aut}_{\FF_p\text{-gr.}}(\SL_{2,\FF_p})\simeq\ad(\SL_{2,\FF_p})$.
\end{enumeratei}
\numpara{7.4.8}\label{XXIV.7.4.8}
It follows in particular that
\begin{enumeratei}
\item $\underline{\Hom}_{\FF_p\text{-gr.}}(\SL_{2,\FF_p},\SL_{2,\FF_p})$ has infinitely many connected components, and therefore \emph{is not quasi-compact}.
\item If $S$ is a scheme of mixed characteristics, $\underline{\Hom}_{S\text{-gr.}}(\SL_{2,S},\SL_{2,S})$ \emph{is not flat} over $S$.
\end{enumeratei}

\sgasection{8}{Appendix: Cohomology of a smooth group over a henselian ring. Cohomology and the functor $\prod$}
\begin{proposition}{8.1}\label{XXIV.8.1}
Let $S$ be a henselian local scheme, $s$ its closed point, and $G$ a smooth $S$-group scheme such that every finite subset of $G$ is contained in an affine open subset.{{\renewcommand{\thefootnote}{*}\addtocounter{footnote}{-1}\footnote{This last condition is in fact unnecessary (cf. A.~Grothendieck, \textit{Groupe de Brauer III}, in \textit{Dix Exposés sur la cohomologie des schémas}, North-Holland, 1968, Theorem 11.7 and Remarks 11.8.3).}}}Then
\begin{enumeratei}
\item If $P$ is a principal homogeneous bundle under $G$, there exists a finite surjective étale morphism $S'\to S$ trivializing $P$. One has $\operatorname{Fib}(S,G)\simeq\mathrm H^1_{\mathrm{\acute et}}(S,G)$.
\item For every finite surjective étale morphism $S'\to S$, the canonical map
\[
\mathrm H^1(S'/S,G)\to\mathrm H^1(S'\otimes_S\kappa(s)/\kappa(s),G_s)
\]
is bijective.
\item The canonical map
\[
\operatorname{Fib}(S,G)\to\operatorname{Fib}(\kappa(s),G_s)
\]
is bijective.
\end{enumeratei}
\end{proposition}
\numpara{8.1.2}\label{XXIV.8.1.2}
If $K$ is a finite separable extension of $\kappa(s)$, there exists a finite surjective étale morphism $S'\to S$ such that $K\simeq S'\otimes_S\kappa(s)$.\nde{Note that $S'$ is still local and henselian. Moreover, we have preserved the numbering of the original: there is no no. 8.1.1.}If $P$ is a principal homogeneous bundle under $G$, then $P$ is smooth over $S$, hence $P_s$ smooth over $\kappa(s)$; there therefore exists a finite separable extension $K$ of $\kappa(s)$ such that $P_K$ has a section (cf. EGA~IV$_4$, 17.15.10). Representing $K$ as indicated above, one sees that $P_{S'}$ has a section by ``Hensel's lemma'' (cf. EGA~IV$_4$, 18.5.17), which proves the first part of (i).

Conversely, if $P$ is a principal homogeneous \emph{sheaf} under $G$ for the étale topology, there exists a finite surjective étale morphism $S'\to S$ trivializing $P$ (indeed every covering family of a henselian local scheme for the étale topology is refined by a family consisting of a single morphism $S'\to S$ of the required form).

By the descent hypothesis made on $G$, $P$ is representable (SGA~1, VIII, 7.6), which completes the proof of (i), and shows that (ii) implies (iii). It therefore only remains for us to prove (ii).
% original PDF p. 50
\numpara{8.1.3}\label{XXIV.8.1.3}
The map of (ii) is injective: let $P$ and $Q$ be two principal homogeneous bundles under $G$ trivialized by $S'$. Consider the $S$-sheaf of groups $H=\underline{\Isom}_{G\text{-bundles}}(P,Q)$; since $H_{S'}$ is isomorphic to $G_{S'}$, $H$ is representable, by the second hypothesis on $G$, cf. above. If $H(\kappa(s))\ne\varnothing$, then $H(S)\ne\varnothing$ by Hensel's lemma, hence $P$ and $Q$ are isomorphic.
% typo?: source calls Isom(P,Q) a sheaf of groups; kept as printed.
\numpara{8.1.4}\label{XXIV.8.1.4}
Let us finally prove that the map of (ii) is surjective, or again that the canonical map
\[
\mathrm Z^1(S'/S,G)\to\mathrm Z^1(S'\otimes_S\kappa(s)/\kappa(s),G_s)
\]
is surjective. Let $\mathbf Z^1$ be the $S$-functor defined by
\[
\mathbf Z^1(T)=\mathrm Z^1(S'\times_S T/T,G_T);
\]
the preceding map identifies with the map
\[
\mathbf Z^1(S)\to\mathbf Z^1(\kappa(s));
\]
by another application of Hensel's lemma, it suffices to prove that $\mathbf Z^1$ is representable by a smooth $S$-scheme.
\numpara{8.1.5}\label{XXIV.8.1.5}
Let us prove that $\mathbf Z^1$ is representable by an $S$-scheme locally of finite presentation. Let $\mathbf C^i$, $i=0,1,\ldots$, be the $S$-functor defined by
\[
\mathbf C^i(T)=\mathrm C^i(S'\times_S T/T,G_T),
\]
i.e.
\[
\mathbf C^i(T)=G((S'\times_S T/T)^{i+1})=G((S'/S)^{i+1}\times_S T),
\]
or again
\[
\mathbf C^i=\underline{\Hom}_S((S'/S)^{i+1},G).
\]
Since $\mathbf Z^1$ is obtained from $\mathbf C^1$ and $\mathbf C^2$ by fibered products, it suffices to prove that $\mathbf C^i$, $i=1,2$, is representable by an $S$-scheme locally of finite presentation.
\numpara{8.1.6}\label{XXIV.8.1.6}
If $T\to S$ is a finite surjective étale morphism splitting $S'$, then
\[
\mathbf C^i\times_S T=\underline{\Hom}_T((S'\times_S T/T)^{i+1},G_T)
\]
is representable by a product of $n$ copies of $G_T$, where $n$ is the degree of $(S'/S)^{i+1}$. Applying the hypothesis on $G$ once more, one deduces that $\mathbf C^i$ is indeed representable by an $S$-scheme locally of finite presentation (SGA~1, VIII, loc. cit.).
\numpara{8.1.7}\label{XXIV.8.1.7}
To prove that $\mathbf Z^1$ is smooth, we must now, by definition, prove that if $T$ is an \emph{affine} $S$-scheme, and $T_0$ the closed subscheme defined by a square-zero ideal $\mathscr J$, the canonical map
\[
\mathbf Z^1(T)\to\mathbf Z^1(T_0)
\]
is surjective. Since $G$ is smooth, the canonical map $G(T)\to G(T_0)$ is surjective, and it suffices to prove that the canonical map
% original PDF p. 51
\[
\mathrm H^1(S'\times_S T/T,G_T)\to\mathrm H^1(S'\times_S T_0/T_0,G_{T_0})
\]
is bijective.

Changing the notation slightly and generalizing the hypotheses, it now suffices for us to prove:
\begin{lemma}{8.1.8}\label{XXIV.8.1.8}
Let $S$ and $S'$ be two \emph{affine} schemes, $S'\to S$ a faithfully flat morphism, $\mathscr J$ a square-zero ideal on $S$, $S_0$ the closed subscheme it defines, and $G$ a smooth $S$-group. Put $S'_0=S'\times_S S_0$, $G_0=G\times_S S_0$. The canonical map
\[
\mathrm H^1(S'/S,G)\to\mathrm H^1(S'_0/S_0,G_0)
\]
is bijective.
\end{lemma}
\begin{remark}{8.1.9}\label{XXIV.8.1.9}
If one assumes $G$ commutative, the same assertion holds for all the $\mathrm H^i$, $i>0$, with an analogous proof.
\end{remark}
\begin{proof}
Let $\mathscr M_0$ be the quasi-coherent $\OO_{S_0}$-module $\mathscr M_0=\Lie(G_0/S_0)\otimes_{\OO_{S_0}}\mathscr J$. For each $S_0$-prescheme $T_0$, put $\mathrm M_0(T_0)=\mathrm H^0(T_0,\mathscr M_0\otimes_{\OO_{S_0}}\OO_{T_0})$, and let
\[
\mathrm M=\prod_{S_0/S}\mathrm M_0\qquad\text{and}\qquad\bar G=\prod_{S_0/S}G_0
\]
be the $S$-functors of groups defined by $\mathrm M(T)=\mathrm M_0(T_0)$ and $\bar G(T)=G_0(T_0)$, where $T_0=T\times_S S_0$. By Exp.~III, 0.9 and (0.6.2), there exists for every affine $S$-scheme $T$ an exact sequence, functorial in $T$:
\[
1\to\mathrm M(T)\to G(T)\to\bar G(T)\to1.
\]
We have to study the map
\[
\mathrm H^1(S'/S,G)\to\mathrm H^1(S'_0/S_0,G_0)=\mathrm H^1(S'/S,\bar G).
\]
Suppose first that $G$ is commutative. One then has an exact cohomology sequence
\[
\mathrm H^1(S'/S,\mathrm M)\to\mathrm H^1(S'/S,G)\to\mathrm H^1(S'/S,\bar G)\to\mathrm H^2(S'/S,\mathrm M);
\]
but
\[
\mathrm H^i(S'/S,\mathrm M)=\mathrm H^i(S'_0/S_0,\mathrm M_0)=\mathrm H^i(S'_0/S_0,\mathscr M_0),
\]
and one knows (TDTE~I, B, Lemma 1.1) that $\mathrm H^i(S'_0/S_0,\mathscr M_0)=0$ for $i\ne0$.

If now $G$ is not commutative, we must use the nonabelian cohomology exact sequence. If $u\in\mathrm Z^1(S'/S,G)$, one knows that the elements of $\mathrm H^1(S'/S,G)$ having the same image in $\mathrm H^1(S'/S,\bar G)$ as the class of $u$ are in the image of the corresponding coboundary map:
\[
\mathrm H^1(S'/S,\mathrm M_u)\to\mathrm H^1(S'/S,G),
\]
where $\mathrm M_u$ is the $S$-functor $\mathrm M$ ``twisted by $u$''. Likewise, if $v$ is an element of $\mathrm Z^1(S'/S,\bar G)$, there exists a ``coboundary''
\[
\Delta(v)\in\mathrm H^2(S'/S,\mathrm M_v),
\]
% original PDF p. 52
where $\mathrm M_v$ is the $S$-functor $\mathrm M$ ``twisted by $v$'', such that $\Delta(v)=0$ if and only if the class of $v$ is in the image of $\mathrm H^1(S'/S,G)$. It suffices for us to prove that $\mathrm H^1(S'/S,\mathrm M_u)=\mathrm H^2(S'/S,\mathrm M_v)=0$.

Now recall (Exp.~III 0.8) that the action of $\bar G$ on $\mathrm M$ defined by the exact sequence is nothing other than the one deduced functorially from the adjoint representation of $G_0$:
\[
\ad:G_0\to\underline{\Aut}_{\OO_{S_0}}(\Lie(G_0/S_0)).
\]
The element $u$ (resp. $v$) therefore acts on $\mathrm M_{S'\times_S S'}$ through an $S'_0\times_{S_0}S'_0$-automorphism of $\Lie(G_0/S_0)$. Since $u$ (resp. $v$) is a cocycle, this automorphism is a descent datum; denote by $\mathscr L_u$ (resp. $\mathscr L_v$) the quasi-coherent $\OO_{S_0}$-module obtained. One verifies at once that for $T\to S$ one has
\[
\mathrm M_u(T)=\mathrm H^0(T_0,\mathscr L_u\otimes_{\OO_{S_0}}\mathscr J\otimes_{\OO_{S_0}}\OO_{T_0})
\]
and the same relation on replacing $u$ by $v$. One therefore has
\[
\begin{aligned}
\mathrm H^1(S'/S,\mathrm M_u)&=\mathrm H^1(S'_0/S_0,\mathscr L_u\otimes\mathscr J),\\
\mathrm H^2(S'/S,\mathrm M_v)&=\mathrm H^2(S'_0/S_0,\mathscr L_v\otimes\mathscr J),
\end{aligned}
\]
and both do indeed vanish by the result already used.
\end{proof}
\begin{proposition}{8.2}\label{XXIV.8.2}
Let $\mathscr C$ be a category possessing fibered products, equipped with a topology coarser than the canonical topology, $S'\to S$ a morphism of $\mathscr C$, $G'$ an $S'$-sheaf of groups, and $G$ the $S$-sheaf of groups $\prod_{S'/S}G'$. Let $\mathrm H^1_S(S',G')\subset\mathrm H^1(S',G')$ be the set of classes of principal homogeneous sheaves under $G'$ which are trivialized by a sieve of $S'$ obtained by base change from a suitable covering sieve of $S$. The canonical map $\mathrm H^1(S,G)\to\mathrm H^1(S',G')$ defined by the functor
\[
P\mto P\times_S S'
\]
induces a bijection
\[
\mathrm H^1(S,G)\isomto\mathrm H^1_S(S',G');
\]
the inverse bijection is defined by the functor $P'\mto\prod_{S'/S}P'$.
\end{proposition}
\begin{proof}
For every object $X$ of $\mathscr C/S$, one has by definition an isomorphism functorial in $X$
\[
\Hom_S(X,G)\isomto\Hom_{S'}(X\times_S S',G').
\]
One therefore has for each $S$-object $T$ a bijection functorial in $T$
\[
\mathrm H^1(T/S,G)\isomto\mathrm H^1(T'/S',G').
\]
Replacing now the single morphism $T\to S$ by an arbitrary covering family of $S$ and passing to the inductive limit, one deduces the first part of the statement. The second part follows without difficulty.
\end{proof}
\begin{lemma}{8.3}\label{XXIV.8.3}
Under the conditions of 8.2, the assertion $\mathrm H^1_S(S',G')=\mathrm H^1(S',G')$ is local on $S$: suppose that there exists a covering family $\{S_i\to S\}$ such that for every $i$ one has $\mathrm H^1_{S_i}(S'\times_S S_i,G')=\mathrm H^1(S'\times_S S_i,G')$. Then $\mathrm H^1_S(S',G')=\mathrm H^1(S',G')$.
\end{lemma}
% original PDF p. 53
\begin{proof}
Indeed, let $P'$ be a principal homogeneous sheaf under $G'$. Put
\[
P'_i=P\times_{S'}(S'\times_S S_i);
\]
% typo?: source has unprimed P on the right; kept as printed.
by the hypothesis, there exists a covering family $\{S_{ij}\to S_i\}$ such that for each $j$, $P'\times_{S'}(S'\times_S S_{ij})$ has a section. But $\{S_{ij}\to S\}$ is a covering family of $S$, and $P'$ is indeed trivialized by the covering family of $S'$ obtained from it by base change.
\end{proof}
\begin{proposition}{8.4}\label{XXIV.8.4}
Let $S'\to S$ be a finite étale morphism of schemes. Let $G'$ be an $S'$-sheaf of groups, and $G$ the $S$-sheaf of groups $\prod_{S'/S}G'$. For the étale (resp. local finite étale, resp. (fpqc)) topology, the functors
\[
\begin{aligned}
P&\mto P\times_S S'\\
\prod_{S'/S}P'&\mathrel{\reflectbox{$\longmapsto$}} P'
\end{aligned}
\]
induce mutually inverse bijections:
\[
\mathrm H^1(S,G)\simeq\mathrm H^1(S',G').
\]
\end{proposition}
\begin{proof}
By 8.2, it suffices to show that $\mathrm H^1_S(S',G')=\mathrm H^1(S',G')$. By 8.3, it suffices to do so locally for the local finite étale topology; one may therefore suppose that $S'$ is a finite direct sum of copies of $S$, say $I_S$, where $I$ is a suitable finite set. Then $G'$ is given by a family $(G_i)_{i\in I}$ of sheaves on $S$, and
\[
\mathrm H^1(S',G')\simeq\prod_{i\in I}\mathrm H^1(S,G_i).
\]
On the other hand
\[
\mathrm H^1(S,G)\simeq\prod_{i\in I}\mathrm H^1(S,G_i),
\]
whence, by 8.2, $\mathrm H^1_S(S',G')=\mathrm H^1(S',G')$.
\end{proof}
\begin{remarks}{8.5}\label{XXIV.8.5}
One can interpret 8.2 and 8.3 by means of the following exact sequence ($f$ is the given morphism $S'\to S$)
\[
1\to\mathrm H^1(S,f_*(G'))\to\mathrm H^1(S',G')\to\mathrm H^0(S,\mathrm R^1f_*(G')).
\]
In the commutative case, this exact sequence follows from the Leray spectral sequence; it is still valid in the noncommutative case (cf. Giraud's thesis\nde{See paragraph V 3.1.4 of [Gi71].}).

In this form, one sees that the result is still valid if $f$ is merely assumed \emph{finite}, or simply \emph{integral}, the topology being the étale topology, since for every $G'$ one has
\[
\mathrm R^1f_*(G')=\text{final sheaf},
\]
by SGA~4, VIII, 5.3.\nde{This reference also points to paragraph V 3.1.4 of [Gi71].}

On the other hand, this result becomes false if one takes a topology such as (fpqc) or (fppf), even if $S=\Spec(k)$, with $k$ an algebraically closed field of characteristic $p\ne0$, $S'=\Spec(k[t]/t^2)$, and $G'=\boldsymbol\mu_p$ or $\boldsymbol\alpha_p$.
% original PDF p. 54

Likewise, 8.2 becomes false, even for the étale topology, if one removes the hypothesis that $f$ is finite, as one sees by taking for $f$ an open immersion; for example, if $S=\Spec(V)$, with $V$ a complete discrete valuation ring with algebraically closed residue field, $S'$ being the open subset induced at the generic point, and $G'$ the constant group $(\ZZ/n\ZZ)_{S'}$, with $n$ prime to the residue characteristic of $V$, one has $\mathrm H^1(S,G)=0$, $\mathrm H^1(S',G')\ne0$. Replacing moreover $S'$ by $S\amalg S'$, one deduces an analogous example, with $S'\to S$ étale and surjective, hence covering for the topology in question.
% typo?: final paragraph cites 8.2, which did not require f finite; kept as printed.
\end{remarks}

\begin{thebibliography}{Jou83}
\bibitem[BLie]{BLie} N.~Bourbaki, \textit{Groupes et algèbres de Lie}, Chap.~I, Hermann, 1960.
\bibitem[Ch51]{Ch51} C.~Chevalley, \textit{Théorie des groupes de Lie. t.~II Groupes algébriques}, Hermann, 1951.
\bibitem[Di57]{Di57} J.~Dieudonné, \textit{Lie groups and Lie hyperalgebras over a field of characteristic $p>0$. VI}, Amer.~J.~Math. \textbf{79} (1957), no.~2, 331--388.\nde{To these three references, appearing in the original, we have added the following references.}
\bibitem[Bo91]{Bo91} A.~Borel, \textit{Linear algebraic groups}, 2nd edition, Springer-Verlag, 1991.
\bibitem[Ch05]{Ch05} C.~Chevalley, \textit{Classification des groupes algébriques semi-simples} (with the collaboration of P.~Cartier, A.~Grothendieck, M.~Lazard), Collected Works, vol.~3, Springer, 2005.
\bibitem[DG70]{DG70} M.~Demazure, P.~Gabriel, \textit{Groupes algébriques}, Masson \& North-Holland, 1970.
\bibitem[Gi71]{Gi71} J.~Giraud, \textit{Cohomologie non abélienne}, Springer-Verlag, 1971.
\bibitem[Ja87]{Ja87} J.~C.~Jantzen, \textit{Representations of Algebraic Groups}, Academic Press, 1987; 2nd edition, Amer.~Math.~Soc., 2003.
\bibitem[Jou83]{Jou83} J.-P.~Jouanolou, \textit{Théorèmes de Bertini et applications}, Birkhäuser, 1983.
% typo: printed Birkhäuser -> Birkhäuser.
\bibitem[Ta75]{Ta75} M.~Takeuchi, \textit{On coverings and hyperalgebras of affine algebraic groups}, Trans.~Amer.~Math.~Soc. \textbf{211} (1975), 249--275.
\end{thebibliography}
